% !TEX root = ../../tree_notes.tex
% =====================================================================
%  Chapter 2 review -- summary, checklist, quick reference, mixed set,
%  related topics. Uses \chaptersection (unnumbered) and NO numbered
%  boxes (they would continue the 2.3.x numbering).
% =====================================================================

\chaptersection{Chapter 2 Summary}

This chapter reduced every binary-tree solution to one recursion and sorted the hard cases
by the \emph{direction information flows}. The tables collect the results and the exact
condition each needs.

\begin{center}
\footnotesize
\renewcommand{\arraystretch}{1.3}
\begin{tabular}{@{}p{0.24\linewidth} p{0.50\linewidth} p{0.20\linewidth}@{}}
\toprule
\textbf{Result} & \textbf{Statement (with conditions)} & \textbf{Where} \\
\midrule
Null links       & A binary tree on $n$ nodes has $n-1$ non-\nil{} child pointers and $n+1$ \nil{} ones. & \thmref{c02-null-links} \\
Bottom-up definedness & With no cycles, the aggregate $g(v)$ depends only on $\subtree(v)$ and always terminates. & \propref{c02-aggregate} \\
Top-down flow    & Answer at $v$ needs ancestor state (depth, running value, grandparent), passed \emph{in}. & \defref{c02-topdown} \\
Bottom-up flow   & Answer at $v$ is a function of the two children's return values, combined and \emph{returned}. & \defref{c02-bottomup} \\
Global-best split& Return the extendable one-side value; record the through-$v$ value in a global. & \coderef{c02-bottomup-tmpl} \\
\bottomrule
\end{tabular}
\end{center}

\medskip
{\small\textbf{Algorithms from this chapter.}}
\begin{center}
\renewcommand{\arraystretch}{1.2}
\begin{tabular}{@{}p{0.36\linewidth} l l l@{}}
\toprule
\textbf{Algorithm} & \textbf{Time} & \textbf{Memory} & \textbf{Code} \\
\midrule
Structural recursion (count / leaves / sum) & $\bigO{n}$ & $\bigO{h}$ & \coderef{c02-kit} \\
Build a tree by linking pointers            & $\bigO{n}$ & $\bigO{n}$ & \coderef{c02-build} \\
Top-down DFS (accumulator / path)           & $\bigO{n}$ (or $\bigO{nh}$ copying) & $\bigO{h}$ & \coderef{c02-topdown-tmpl} \\
Bottom-up DFS (aggregate + global best)     & $\bigO{n}$ & $\bigO{h}$ & \coderef{c02-bottomup-tmpl} \\
\bottomrule
\end{tabular}
\end{center}

\begin{figure}[H]
\centering
\begin{tikzpicture}[diagram, node distance=7mm,
    q/.style={boxnode, fill=primaryblue!12, draw=primaryblue, text width=32mm},
    a/.style={boxnode, fill=successgreen!12, draw=successgreen, text width=34mm}]
  \node[q] (ask) {\textbf{At a node, what does the answer need?}};
  \node[a, below left=10mm and 0mm of ask] (m) {just the whole subtree \\ \footnotesize count / collect / compare (\secref{c02-model})};
  \node[a, below=10mm of ask] (d) {something about the \emph{ancestors} \\ \footnotesize depth, running sum, grandparent $\Rightarrow$ top-down (\secref{c02-topdown})};
  \node[a, below right=10mm and 0mm of ask] (u) {something about the \emph{descendants} \\ \footnotesize subtree sum/height/best $\Rightarrow$ bottom-up (\secref{c02-bottomup})};
  \draw[flow] (ask) -- (m); \draw[flow] (ask) -- (d); \draw[flow] (ask) -- (u);
\end{tikzpicture}
\caption*{Choosing the recursion by data direction. If a node's answer needs the path above
it, carry state down (2.2); if it needs the subtrees below, return an aggregate up (2.3);
if it just needs to sweep the whole tree, plain structural recursion (2.1) suffices.}
\end{figure}

\chaptersection{Common Mistakes Checklist}

\begin{itemize}[nosep, leftmargin=1.4em]
  \item \textbf{No \nil{} base case.} You meet $n+1$ nulls (\thmref{c02-null-links}); dereferencing
        one crashes. Fix: test \code{r == nullptr} first, every function.
  \item \textbf{Leaf vs null confused.} A leaf is a real node with two null children; a null
        is the absence of a node. Fix: give each its own return value.
  \item \textbf{Shared state not restored.} A pushed path or incremented counter leaks to the
        next sibling. Fix: undo after recursing, or pass primitives by value (\probref{c02-root-numbers}).
  \item \textbf{Returning the through-node answer.} Returning the bend-at-$v$ value makes the
        parent double-count. Fix: return one side, keep the through value in a global.
  \item \textbf{Wrong identity at null.} Sums want $0$, edge-heights want $-1$, prune wants
        \code{nullptr}. Fix: pick the identity from the combine $\phi$, then test a one-node tree.
  \item \textbf{Edges vs nodes again.} Univalue/diameter count edges; depth counts nodes.
        Same off-by-one as \secref{c01-anatomy}.
  \item \textbf{String / vector concatenation by value.} \probref{c02-paths} and \probref{c02-construct}
        become $\bigO{n h}$ or $\bigO{n^2}$. Fix: append into one shared buffer.
\end{itemize}

\chaptersection{Quick Reference Card}

\begin{quickref}[Chapter 2 at a glance]
\textbf{The node.} \code{TreeNode\{val,left,right\}}; empty tree $=$ \code{nullptr}; two
distinguished children. $n$ nodes $\Rightarrow$ $n-1$ live links, $n+1$ null links.

\medskip
\textbf{Structural recursion (2.1).} Base at \nil; combine two correct subtree answers with
the node. Count/leaves/sum/paths/serialise all share this one skeleton.

\medskip
\textbf{Top-down (2.2).} $f(v,\text{state})$; state is ancestor info (depth, accumulator,
path, grandparent) refined on the way \emph{in}; restore shared state between siblings.

\medskip
\textbf{Bottom-up (2.3).} Return $g(v)=\phi(v,g_L,g_R)$ \emph{out}; when the answer passes
through $v$ but the parent needs one side, \emph{return} the side, \emph{record} the through
value in a global (tilt, univalue path, diameter).
\end{quickref}

\begin{lstlisting}[caption={Chapter 2 code kit: the two recursion directions side by side}, label={lst:c02-kit2}]
#include <bits/stdc++.h>
using namespace std;
struct TreeNode { int val; TreeNode *left, *right;
                  TreeNode(int x=0): val(x), left(nullptr), right(nullptr){} };

// TOP-DOWN: carry a value down; record at leaves. (Sum Root to Leaf Numbers, LC 129)
int topDown(TreeNode* r, int acc) {
    if (!r) return 0;
    acc = acc * 10 + r->val;                          // refine state for children
    if (!r->left && !r->right) return acc;            // leaf: emit
    return topDown(r->left, acc) + topDown(r->right, acc);
}

// BOTTOM-UP: return an aggregate, update a global best. (Longest Univalue Path, LC 687)
int best = 0;
int chain(TreeNode* r) {
    if (!r) return 0;
    int L = chain(r->left), R = chain(r->right);
    int le = (r->left  && r->left->val  == r->val) ? L + 1 : 0;   // extendable
    int re = (r->right && r->right->val == r->val) ? R + 1 : 0;
    best = max(best, le + re);                        // answer THROUGH r
    return max(le, re);                               // return ONE side up
}
\end{lstlisting}

\chaptersection{Mixed Problem Set}

These blend the chapter's ideas and never say which recursion direction to use; decide it
yourself from what each node's answer depends on.

\begin{mixedset}
  \practice{Find Bottom Left Tree Value}{LeetCode 513}
    {https://leetcode.com/problems/find-bottom-left-tree-value/}{LC Medium}
    {carry depth down and keep the first node seen at each new maximum depth.}
  \practice{Add One Row to Tree}{LeetCode 623}
    {https://leetcode.com/problems/add-one-row-to-tree/}{LC Medium}
    {top-down to depth $d-1$, then splice in two new children; special-case $d=1$.}
  \practice{Find Duplicate Subtrees}{LeetCode 652}
    {https://leetcode.com/problems/find-duplicate-subtrees/}{LC Medium}
    {bottom-up: serialise each subtree into a string, hash it, and report hashes seen twice.}
  \practice{Flatten Binary Tree to Linked List}{LeetCode 114}
    {https://leetcode.com/problems/flatten-binary-tree-to-linked-list/}{LC Medium}
    {pre-order, but rewire right pointers; the in-place trick splices the old right subtree.}
  \practice{Flip Equivalent Binary Trees}{LeetCode 951}
    {https://leetcode.com/problems/flip-equivalent-binary-trees/}{LC Medium}
    {same or flipped: recurse comparing (a,b) and (b,a) child pairs at every node.}
  \practicehh{Binary Tree Cameras}{LeetCode 968}
    {https://leetcode.com/problems/binary-tree-cameras/}{LC Hard}
    {Stretch: greedy bottom-up -- return each node's coverage state to its parent.}
    {three states (has camera / covered / needs coverage); decide a camera when a child needs coverage (\chref{c10} reroot-free DP).}
\end{mixedset}

\chaptersection{Related Topics}

\begin{itemize}[nosep, leftmargin=1.4em]
  \item \chref{c01} -- the tree vocabulary and the height/size recursion this chapter generalises.
  \item \chref{c03} -- preorder/inorder/postorder/level-order: the walks these recursions embody.
  \item \chref{c04} -- array, adjacency-list and Euler-tour storage; how to \emph{hold} the tree.
  \item \chref{c05} -- binary search trees, where inorder order adds new bottom-up tricks.
  \item \chref{c09} -- diameter and centroid: the return-plus-global idiom on general trees.
  \item \chref{c10} -- tree DP formalises the bottom-up aggregate with richer returned states.
  \item \chref{c11} -- LCA: carrying parent/depth information down, the ancestor tool teased here.
\end{itemize}
